\chapter{Workspace Architecture}

\section{Repository Layout}

The workspace keeps root-level scripts as user-facing entry points and stores
implementation details in importable modules under \code{src/}. Source data,
processed data, reports, and figures are separated to make outputs easy to
inspect and clean.

\begin{table}[htbp]
\centering
\caption{Core workspace directories.}
\begin{tabular}{@{}p{0.28\linewidth}p{0.58\linewidth}@{}}
\toprule
Path & Role \\
\midrule
\code{src/} & Reusable Python modules. \\
\code{src/tools/} & Data parsing, TE plotting, XRD utilities, and SPB fitting. \\
\code{data/raw/} & Original instrument files and raw batch folders. \\
\code{data/processed/} & Cleaned transport-property CSV tables. \\
\code{data/lab/} & Batch/sample ledgers and editable lab metadata. \\
\code{results/} & Reports, assessment payloads, fitting outputs, predictions. \\
\code{outputs/} & Figures, one-off exports, tables, and logs. \\
\bottomrule
\end{tabular}
\end{table}

\section{Planned Content}

\begin{itemize}[leftmargin=*]
  \item Data ownership rules for raw, processed, and generated files.
  \item Naming conventions for CHY batches and samples.
  \item Relationship between \code{results/} and \code{outputs/}.
  \item How external snapshots under \code{external/} should be treated.
\end{itemize}
